\section{確率分布}
\label{lesson06}

\subsection{授業内容}
\subsubsection{科目の中でのこのコマの位置づけ}
\begin{tabular}[t]{ccccccc}
    記述統計[3] & 確率[1] & モデル[1] & 推測・推定[0] & 意思決定[0] & 実践[1]
\end{tabular}
\subsubsection{概要}
ここでは統計学と密接なつながりを持つ，確率と確率分布について学ぶ。記述統計ではデータの中心的な位置，ばらつきの幅を記述することができたが，それらがどのような規則性を持つかを記述するのが確率分布の考え方である。確率の基本的な性質と，生じうる結果の一覧としての確率分布を理解する。代表例としてここでは二項分布を扱う。これらは記述統計の枠を超えた推測統計，確率モデルへの基礎的知識となる。
\subsubsection{コマ主題細目}
\begin{description}

\item[確率と呼ばれる数字の二つの理解]高校数学までの範囲で学ぶ確率は，全ての場合の数を全体としたうち，該当する事象の生じる回数を相対頻度で表すものであった。一方我々は，天気予報で降水確率が30\%というように，確率的な表現を日常的にも用いる。天気予報は未来の予測であるから，「起こりうる全ての未来」を書き出して生じる天気の割合を表現できるはずもない。ここには二種類の確率の考え方が含まれているし，同時に両者とも共有しているルールがある。
\begin{flushright}
    $\to$ \citet{Maeda2017}のP.78--83
\end{flushright}

\item[確率分布]確率的な事象と見立てて確率の言葉で記述するとき，起こりうる全ての可能性を網羅的に表したものを確率分布という。離散変数の場合は確率質量関数とよばれる，全ての起こりうるカテゴリに付与された確率の一覧がその確率分布となる。連続的な分布の場合は確率密度関数と呼ばれる関数に変わることに注意が必要である。連続変数の場合，一点の確率は定義できず，区間の面積が確率の数字になる。
\begin{flushright}
    $\to$ \citet{Maeda2017}のP.83--88
\end{flushright}

\item[二項分布]確率分布の中でも最も単純なのがベルヌーイ分布であり，それを拡張したのが二項分布である。離散変数の代表的な例としてベルヌーイ分布と二項分布について，数式や表記法とともに理解し，また確率分布の平均値である期待値，確率分布の分散を表す方法を学ぶ。
\begin{flushright}
    $\to$ \citet{野島201903}のP.66-68
\end{flushright}

\end{description}
\subsubsection{キーワード}
\begin{itemize}
\item 確率
\item 頻度による定義と主観による定義
\item ベイズ統計
\item 確率分布
\item ベルヌーイ試行，ベルヌーイ分布，二項分布
\end{itemize}
\subsection{授業情報}
\paragraph{コマの展開方法}
講義
\paragraph{標準シラバスにおける位置づけ}
科目番号5 心理学統計法；1.心理学で用いられる統計手法；D 離散分布：二項分布を中心として
\subsubsection{予習・復習課題}
\paragraph{予習}
これまで習ってきた高校数学の，確率・統計に関する記述を改めて読み直してこの授業への準備としよう。確率についての数学は比較的歴史が浅く，また専門家でも間違えることの多い学問領域であるが，同時にギャンブルや保険の利率計算など生活に密着した数字でもある。副読本として読み物\citet{浜田201812}を挙げておくので，最初の数章だけでも目を通しておくと心構えができるだろう。
\paragraph{復習}
確率分布は記述統計学から推測統計学へと展開する際の，最初の難関であると思われる。高校数学までの「相対頻度」によって表されるのとは異なる「主観確率」の考え方がある。前者を頻度主義，後者をベイズ主義(ベイジアン)と呼んで区別することもあるが，確率に関するこの区別は主義・主張の派閥の対立というより，説明のための言葉使いの問題であることに留意し，そもそも確率という数字が従うべき共通のルールをしっかり把握しておくと良い。\citet{Maeda2017}のCh.1は主観確率を積極的に用いるベイズ統計学の入門として非常にわかりやすい例が含まれている。
